\documentclass[10pt, a4paper, twocolumn]{article}

\usepackage{fontspec}
\usepackage{polyglossia}
\setdefaultlanguage{romanian}
\setotherlanguages{english}
\SetLanguageKeys{romanian}{indentfirst=true}

\usepackage{csquotes}

\usepackage[
    top=2cm, 
    bottom=2cm, 
    left=1.5cm, 
    right=1.5cm, 
    columnsep=0.6cm
]{geometry}

\usepackage{setspace}
\singlespacing 

\usepackage{graphicx}
\usepackage{float}
\usepackage{booktabs}
\usepackage{array}
\usepackage{enumitem}

\usepackage{caption}
\usepackage{subcaption}

\usepackage{amsmath}
\usepackage{mathtools}
\usepackage{algorithm}
\usepackage{algpseudocode}
\usepackage[warnings-off={mathtools-colon,mathtools-overbracket}]{unicode-math}
\newcommand{\bigO}[1]{\symcal{O}\left(#1\right)}
\DeclarePairedDelimiter\abs{\lvert}{\rvert}

\usepackage{listings}
\usepackage{xcolor}

\lstset{
    language=Python,
    basicstyle=\ttfamily\scriptsize,
    keywordstyle=\color{blue}\bfseries,
    commentstyle=\color{green!40!black},
    stringstyle=\color{orange},
    numbers=left,
    numberstyle=\tiny\color{gray},
    stepnumber=1,
    numbersep=3pt,
    breaklines=true,
    breakatwhitespace=true,
    frame=single,
    backgroundcolor=\color{gray!5},
    captionpos=b,
    tabsize=2,
    showspaces=false,
    showstringspaces=false
}

\usepackage[
    maxbibnames=50,
    sorting=nty,
    backend=biber
]{biblatex}
\addbibresource{bibliography.bib}

\usepackage[
    colorlinks,
    linkcolor={black},
    citecolor={blue},
    urlcolor={blue}
]{hyperref}

\usepackage{titling}
\pretitle{\begin{center}\LARGE\bfseries}
\posttitle{\par\end{center}\vskip 0.5em}
\preauthor{\begin{center}\normalsize}
\postauthor{\end{center}}
\predate{\begin{center}\normalsize}
\postdate{\end{center}}

\usepackage{fancyhdr}
\pagestyle{fancy}
\fancyhf{}
\fancyhead[L]{\small \textit{Algoritmi Paraleli și Distribuiți}}
\fancyhead[R]{\small \textit{APD 2026}}
\fancyfoot[C]{\thepage}
\renewcommand{\headrulewidth}{0.5pt}

% Titlul „Rezumat” centrat în coloană
\makeatletter
\renewenvironment{abstract}{%
  \if@twocolumn
    \begin{center}\textbf{\abstractname}\end{center}\par\vspace{0.5em}%
  \else
    \small
    \begin{center}%
      \bfseries \abstractname\vspace{-.5em}\vspace{\z@}%
    \end{center}%
    \quotation
  \fi
}{%
  \if@twocolumn
  \else
    \endquotation
  \fi
}
\makeatother

\title{\textbf{Monitorizare Distribuită prin Protocolul Gossip}\\
\large Analiză și Implementare a Protocolului Simetric Push-Pull}

\author{
    \href{https://www.linkedin.com/in/andrei-murariu-83317b177/}{\textbf{Murariu Andrei-Stefanel}} \\
    \textit{Facultatea de Matematică și Informatică, Universitatea din București} \\
    \texttt{andrei-stefanel.murariu@s.unibuc.ro}
}
\date{\today}

\begin{document}

\twocolumn[
  \begin{@twocolumnfalse}
    \maketitle
    \begin{abstract}
        Această lucrare analizează eficiența protocoalelor de tip Gossip pentru monitorizarea sistemelor distribuite la scară largă, propunând o alternativă descentralizată la arhitecturile clasice Client-Server. Se investighează algoritmul \textit{Push-Pull Dynamic Average Consensus}, evaluând performanța acestuia pe diverse topologii de rețea (Mesh, Ring, Star). Rezultatele experimentale demonstrează că abordarea distribuită elimină punctul unic de eșec (SPOF) și oferă o scalabilitate superioară. De asemenea, lucrarea tratează fenomenul de \textit{Mass Loss} în scenarii de eșec al nodurilor și impactul distribuției inegale a valorilor asupra timpului de convergență.
    \end{abstract}
    \vspace{1em}
  \end{@twocolumnfalse}
]
\section{Introducere}
\label{sec:intro}

În contextul tehnologic actual, caracterizat prin expansiunea accelerată a paradigmei IoT (\textit{Internet of Things}), arhitecturile sistemelor distribuite au suferit transformări fundamentale. Odată cu migrarea masivă a serviciilor către infrastructuri de tip Cloud-native, fie sub formă de SaaS (\textit{Software as a Service}), PaaS sau arhitecturi Serverless, complexitatea operațională a crescut exponențial. Aceste medii dinamice impun cerințe stricte privind observabilitatea, monitorizarea și controlul în timp real al resurselor.

Problematica monitorizării infrastructurilor moderne provine din limitările inherente ale arhitecturilor centralizate tradiționale. Deși eficiente la scară mică, sistemele bazate pe un coordonator unic care interoghează agenți pasivi devin rapid ineficiente în fața volumului de date generat de mii de containere sau microservicii efemere. Arhitecturile centralizate prezintă două vulnerabilități critice:

\begin{itemize}
    \item \textbf{Dependența de un Punct Unic de Eșec (SPOF):} În modelul clasic, starea globală a sistemului este agregată de un singur nod central. Această centralizare creează un \textit{Single Point of Failure}. Indisponibilitatea serverului central, fie din cauze hardware, fie software, duce la pierderea totală a vizibilității asupra întregii infrastructuri, lăsând administratorii fără date critice exact în momentele de criză.
    
    \item \textbf{Limitări de Scalabilitate și Latență:} Pe măsură ce numărul de noduri monitorizate ($N$) crește, pe nodul central întâmplă un proces de \textit{bottleneck}. Procesarea secvențială sau concurentă a mii de cereri de stare duce la saturarea lățimii de bandă și a resurselor CPU. Soluțiile convenționale implică costuri ridicate:
    \begin{itemize}
        \item \textit{Scalabilitate verticală:} Upgrade-ul hardware al serverului central este o soluție temporară și extrem de costisitoare.
        \item \textit{Scalabilitate orizontală:} Introducerea de proxy-uri sau sharding adaugă un nivel de complexitate în configurare și mentenanță care poate depăși beneficiile.
    \end{itemize}
\end{itemize}

Această lucrare propune și analizează o soluție alternativă, robustă și descentralizată: utilizarea unui protocol de tip \textbf{Gossip} bazat pe strategia \textbf{Dynamic Average Consensus} \cite{spanos2005dynamic}. Prin adoptarea unui model de comunicare simetric \textit{Push-Pull}, se demonstrează posibilitatea monitorizării în timp real a stării globale a sistemului (ex. încărcarea medie a procesorului) fără a utiliza un coordonator central, eliminând astfel riscurile structurale ale abordărilor clasice.
\section{Fundamente Teoretice}
\label{sec:background}

\subsection{Protocoale de tip Gossip}
Protocoalele de tip Gossip, modelează diseminarea informației într-un sistem distribuit similar modului în care se răspândește un zvon într-o rețea socială. Caracteristica definitorie a acestor protocoale este absența unei autorități centrale care să coordoneze comunicarea; în schimb, informația se propagă prin schimburi periodice între perechi de noduri.

Conform taxonomiei propuse de Shah \cite{shah2009gossip}, eficiența unui algoritm Gossip într-un mediu distribuit este condiționată de respectarea a patru principii fundamentale:
\begin{itemize} 
    \item \textbf{Localitate Strictă:} Un nod deține informații exclusiv despre vecinii săi direcți din graful de rețea. Nu este necesară menținerea unei tabele globale de rutare sau cunoașterea topologiei complete a rețelei.
    \item \textbf{Simplitate Computațională:} Algoritmul trebuie să implice operații aritmetice elementare și un consum minim de memorie, permițând rularea pe dispozitive cu resurse limitate. Interacțiunea nu trebuie să implice transferuri masive de date.
    \item \textbf{Asincronicitate:} Sistemul nu presupune existența unui ceas global sau a unor runde de comunicare perfect sincronizate. Nodurile inițiază schimburi de date independent, bazându-se pe ceasuri locale.
    \item \textbf{Toleranță la Erori (Fault Tolerance):} Algoritmul trebuie să fie robust la schimbări dinamice ale topologiei, cum ar fi defectarea nodurilor sau pierderea pachetelor, asigurând convergența rezultatului final fără intervenție externă.
\end{itemize}

\subsection{Consensul Distribuit (Average Consensus)}
Problema fundamentală abordată în această lucrare este \textit{Average Consensus}: convergența stării tuturor nodurilor către media aritmetică a valorilor inițiale din rețea. Fie $x_i(0)$ valoarea inițială a nodului $i$ (de exemplu, încărcarea CPU). Obiectivul este ca, printr-un proces iterativ, starea fiecărui nod $x_i(t)$ să satisfacă condiția \cite{spanos2005dynamic}:
\[ \lim_{t \to \infty} x_i(t) = \frac{1}{n} \sum_{j=1}^{n} x_j(0) \]

\subsection{Uniform Gossip și Conservarea Masei}
Modelul utilizat, \textbf{Uniform Gossip} \cite{kempe2003gossip}, presupune că fiecare nod selectează un partener de comunicare cu o probabilitate uniformă $1/n$ (în cazul unei topologii complet conectate). Mecanismul de actualizare se bazează pe proprietatea de \textbf{Conservare a Masei}.

Atunci când două noduri, $i$ și $j$, interacționează la momentul $t$, ele își actualizează stările la media valorilor lor curente:
\[ v_i(t+1) = v_j(t+1) = \frac{v_i(t) + v_j(t)}{2} \]
Această operație asigură invarianța sumei globale:
\[ v_i(t+1) + v_j(t+1) = v_i(t) + v_j(t) \]
Astfel, informația nu este pierdută sau creată, ci doar redistribuită, conducând sistemul către o stare de echilibru (entropie maximă).

\subsection{Dynamic Consensus}
În scenarii reale, valorile monitorizate nu sunt statice. Deoarece reinițializarea algoritmului la fiecare fluctuație a senzorilor este impracticabilă, se utilizează \textit{Dynamic Average Consensus}. Nodurile monitorizează continuu valoarea locală și, în loc să suprascrie starea algoritmului, "injectează" diferența ($\Delta$) dintre noua măsurătoare și cea anterioară direct în procesul de mediere. Aceasta permite urmărirea în timp real a mediei globale variabile \cite{spanos2005dynamic}.

\subsection{Impactul Topologiei de Rețea}
Performanța algoritmilor de consens este dictată de spectrul grafului de conexiuni:
\begin{itemize} 
    \item \textbf{Topologia Inel:} Fiecare nod are gradul $k=2$. Deși simplă, prezintă o vulnerabilitate majoră la partiționare și o viteză de convergență redusă $\bigO{N^2}$.
    \item \textbf{Topologia Stea:} Centralizează comunicarea într-un Hub. Deși diametrul este mic, Hub-ul reintroduce problema SPOF și a congestiei, contrazicând principiile distribuite.
    \item \textbf{Topologia Completă:} Permite comunicarea directă între oricare două noduri. Aceasta minimizează diametrul rețelei și maximizează viteza de amestecare, pentru că are convergență rapidă ($\bigO{\log N}$).
\end{itemize}
\section{Implementare și Algoritmi}
\label{sec:algorithms}

Pentru a valida ipotezele teoretice, am implementat și comparat două arhitecturi distincte: un model serial de referință și algoritmul distribuit propus.

\subsection{Algoritmul Serial de Referință}
Ca punct de referință, am modelat abordarea centralizată standard. Un server central acționează ca un coordonator care interoghează secvențial sau concurent agenții de monitorizare.

Fluxul operațional este structurat în cicluri de monitorizare:
\begin{enumerate}
    \item \textbf{Colectarea (Polling):} Serverul recepționează valorile $v_i$ de la toate cele $N$ noduri active.
    \item \textbf{Calculul:} Serverul calculează suma totală $S_{total} = \sum v_i$ și derivă media aritmetică $\mu = S_{total} / N$.
\end{enumerate}

Deși acest algoritm oferă o precizie matematică absolută (în absența erorilor de rețea), el scalează liniar $\bigO{N}$ în ceea ce privește timpul de procesare și resursele necesare, devenind impracticabil în rețelele masive din cauza limitărilor de lățime de bandă la nivelul serverului \cite{van2017distributed}.

\subsection{Algoritmul Distribuit: Push-Pull Gossip}
Soluția propusă adoptă modelul \textit{SPMD (Single Program, Multiple Data)}, unde toate nodurile sunt echivalente funcțional. Arhitectura elimină coordonatorul central, bazându-se pe fire de execuție independente pentru emisia și recepția datelor.

Algoritmul utilizează protocolul \textbf{Push-Pull}, care este superior variantelor \textit{Push-Only} deoarece reduce timpul de convergență prin actualizarea simultană a stării ambelor noduri participante la o interacțiune. Pseudocodul \ref{alg:push-pull} descrie logica implementată pe fiecare nod.

\begin{algorithm}[H]
    \caption{Push-Pull Dynamic Average Consensus}
    \label{alg:push-pull}
    \begin{algorithmic}[1]
        \State Inițializează lista de vecini conform topologiei
        \State \textbf{Start Paralel: Thread Activ, Thread Pasiv}

        \State \textbf{Thread 1: Activ}
        \State \quad Măsoară încărcarea CPU curentă ($L_{curr}$)
        \State \quad \textbf{Atomic:} Actualizează media cu diferența $(L_{curr} - L_{last})$
        \State \quad Selectează aleatoriu un vecin $V$
        \State \quad Trimite media locală către $V$ (PUSH)
        \State \quad Primește media vecinului $V$ (PULL)
        \State \quad \textbf{Atomic:} $Media \leftarrow (Media + Media_{vecin}) / 2$

        \State \textbf{Thread 2: Pasiv}
        \State \quad Așteaptă conexiune de la un inițiator
        \State \quad Primește valoarea $Val_{in}$
        \State \quad \quad \textbf{Atomic:} Salvează media curentă în $R$
        \State \quad \quad \textbf{Atomic:} $Media \leftarrow (Media + Val_{in}) / 2$
        \State \quad Trimite $R$ înapoi către inițiator
    \end{algorithmic}
\end{algorithm}

\vspace{0.5em}
\noindent \textbf{Notă:} Operațiile \textit{Atomic} sunt protejate de un \textbf{Mutex} pentru a evita conflictele. Astfel, ne asigurăm că firele de execuție nu scriu simultan în aceeași variabilă, fapt ce ar duce la pierderea datelor și la un calcul greșit al mediei.

\subsection{Mecanismul de Eșec: "Negative Mass Injection"}
Un aspect tehnic crucial observat în implementarea \textit{Dynamic Consensus} este comportamentul sistemului la ieșirea bruscă a unui nod din rețea. Algoritmul se bazează pe injecția diferențelor ($\Delta$). 

Dacă un nod $A$ a partajat deja o valoare mare (ex. 100\%) cu rețeaua, "masa" sa a fost transferată parțial către vecini. Dacă nodul $A$ părăsește rețeaua (crash), el ia cu el valoarea sa curentă. Totuși, din perspectiva algoritmului distribuit, dacă nodul $A$ ar reveni cu valoarea 0, el ar injecta un $\Delta = 0 - 100 = -100$. 

În testele noastre de robustețe, am observat că dispariția nodurilor fără un protocol de \textit{graceful shutdown} poate duce la pierderea definitivă a unei cantități din suma globală ("Mass Loss"). Acest fenomen explică de ce, în anumite scenarii de avarie, media estimată a rețelei scade brusc, sistemul stabilizându-se la o nouă valoare de consens care reflectă "masa" rămasă în nodurile active \cite{kempe2003gossip}.
\section{Analiză Teoretică și Discuții}
\label{sec:analysis}

În această secțiune, analizăm performanța algoritmului distribuit propus (Push-Pull Dynamic Average Consensus) comparativ cu soluția serială, urmărind criteriile de eficiență, scalabilitate și impactul topologiei.

\subsection{Complexitate Timp și Mesaje}
Unul dintre avantajele majore ale algoritmilor Gossip este reducerea complexității temporale în rețelele mari.

\begin{itemize}
    \item \textbf{Convergența:}
    \begin{itemize}
        \item \textbf{Algoritmul Serial:} Necesită $O(N)$ timp per ciclu, deoarece serverul central trebuie să proceseze secvențial $N$ mesaje. Timpul crește liniar cu numărul de noduri.
        \item \textbf{Algoritmul Gossip:} Conform teoriei demonstrate de Kempe \cite{kempe2003gossip}, într-o topologie completă (sau random graph), timpul necesar pentru ca eroarea estimării să scadă la jumătate este de ordinul $O(\log N)$. Aceasta înseamnă o convergență exponențială: dacă dublăm numărul de noduri, timpul de convergență crește doar cu o valoare constantă mică.
    \end{itemize}
    
    \item \textbf{Mesaje:}
    \begin{itemize}
        \item În abordarea \textbf{Serială}, serverul central gestionează $2N$ mesaje pe ciclu (Request + Reply pentru fiecare nod), ceea ce duce la congestie.
        \item În abordarea \textbf{Gossip}, fiecare nod trimite un număr constant de mesaje per unitatea de timp (de obicei 1 mesaj către 1 vecin). Deși numărul total de mesaje în rețea este $O(N)$, acestea sunt distribuite uniform pe toate link-urile, evitând gâtuirea rețelei.
    \end{itemize}
\end{itemize}

\subsection{Speedup și Eficiență}
Accelerația (Speedup) măsoară cât de mult câștigăm prin paralelizare. Deoarece algoritmul Gossip rulează simultan pe toate cele $N$ noduri (model SPMD), capacitatea totală de procesare a sistemului crește odată cu dimensiunea rețelei.

\begin{itemize}
    \item \textbf{Utilizarea Resurselor:} Spre deosebire de modelul Client-Server unde serverul stă la 100\% CPU iar clienții la 1\% (așteptând), în Gossip efortul este partajat egal. Eficiența este maximă deoarece nu există timpi morți cauzați de așteptarea la coadă la un coordonator central.
    \item \textbf{Reducerea Erorii:} Eficiența algoritmului este vizibilă în viteza cu care media estimată se apropie de media reală. În testele efectuate, eroarea scade cu un factor constant la fiecare rundă de comunicare ($e^{-t}$), fenomen numit "convergență geometrică".
\end{itemize}

\subsection{Avantaje și Dezavantaje la Variația Numărului de Noduri}

\subsection{Particularizare pe Topologii}
Performanța algoritmului este dictată strict de structura grafului de conexiuni ($G=(V,E)$). În cadrul proiectului, am analizat două cazuri extreme:

\textbf{1. Topologia Inel (Ring):}
\begin{itemize}
    \item \textit{Descriere:} Fiecare nod comunică doar cu vecinii stânga/dreapta ($grad=2$).
    \item \textit{Performanță:} Slabă. Convergența este lentă, de ordinul $O(N^2)$. Informația trebuie să parcurgă fizic tot lanțul pentru a ajunge de la un capăt la altul.
\end{itemize}

\textbf{2. Topologia Completă (Mesh):}
\begin{itemize}
    \item \textit{Descriere:} Orice nod poate contacta orice alt nod.
    \item \textit{Performanță:} Excelentă. Convergența este de ordinul $O(\log N)$. Diametrul mic al grafului permite informației să se propage "viral" în câțiva pași.
\end{itemize}
\section{Concluzii}
\label{sec:conclusion}

Lucrarea de față a analizat viabilitatea utilizării protocolului \textit{Push-Pull Dynamic Average Consensus} pentru monitorizarea sistemelor distribuite, ca alternativă la arhitecturile centralizate. Rezultatele obținute confirmă ipoteza că abordarea descentralizată oferă avantaje majore în termeni de scalabilitate și reziliență.

Prin comparația directă a topologiilor, am demonstrat că arhitectura \textbf{Mesh} oferă cel mai bun compromis pentru convergența rapidă ($\bigO{\log N}$), fiind superioară topologiilor Ring (prea lentă) și Star (vulnerabilă la congestie). De asemenea, testele de robustețe au evidențiat capacitatea algoritmului de a se adapta dinamic la schimbările de infrastructură. Deși pierderea bruscă a nodurilor cauzează o fluctuație a mediei estimate (\textit{Mass Loss}), sistemul nu intră în colaps, ci converge către o nouă stare stabilă. 
\printbibliography[heading=bibintoc]

\end{document}